\begin{exo}{}
	On considère une fonction $f$ définie sur $\intff{-3}{3}$ dont un tableau de valeurs est donné ci-dessous :


	\begin{center}
		\renewcommand{\arraystretch}{1.5}
		\begin{tabular}{|c|C{5mm}|C{5mm}|C{5mm}|C{5mm}|C{5mm}|C{5mm}|C{5mm}|} \hline
			\cellcolor{gray!30}$x$    & $-3$ & $-2$ & $-1$ & $0$ & $1$ & $2$ & $3$ \\ \hline
			\cellcolor{gray!30}$f(x)$ & $4$  & $1$  & $0$  & $1$ & $2$ & $3$ & $2$ \\ \hline
		\end{tabular}
	\end{center}
	\begin{enumerate}
		\item Des antécédents de $2$ par $f$ sont :
		      \vspace*{-2mm}
		      \begin{multicols}{2}
		      \begin{enumerate}
			      \item $1$
			      \item $2$
			      \item $3$
			      \item Pas d'antécédent
		      \end{enumerate}
		      \end{multicols}
		      \vspace*{-4mm}
		\item Des antécédents de $3$ par $f$ sont :
		      \begin{enumerate}
			      \item $2$
			      \item $3$
			      \item On ne peut pas connaître tous les antécédents.
			      \item $3$ n'a pas d'antécédent.
		      \end{enumerate}
	\end{enumerate}
\end{exo}
